\documentclass[%
        paper=a4,
	% BCOR=10mm, % Binding CORrection
	oneside,			% two sided printing
        %cleardoublepage=empty,         % empty backsides for two sided printing
        twocolumn=false,
	titlepage=false,			% support for title environment
        abstract,
	numbers=noenddot,		% numbers (in headlines) without ending points
	headsepline,			% line below header
	%footsepline,			% line above footer
	headings=normal,                % size of header is normal
	fontsize=11pt,			% fontsize in points
	pagesize=pdftex,		% define typearea for pdftex usage
        headheight=46.96317pt,          % value based on MyThOS logo
        parskip=off                     % or half*
]{scrartcl}

% ------------------------------------------------------------------------------
% character encoding and fonts for the document
\usepackage[T1]{fontenc} % should be loaded before inputenc
\usepackage[utf8]{inputenc}	    % unicode, newer than utf8x
\usepackage{ae}
\usepackage{aecompl}
\usepackage[T1]{url}			% web adresses with t1 style
\urlstyle{tt}					% web adresses in tt-style
%\usepackage[UKenglish]{babel}
\usepackage[ngerman,UKenglish]{babel}

% ------------------------------------------------------------------------------
% font setup
\usepackage{microtype}
\usepackage{ellipsis}			% corrects wrong spaces after dots
\usepackage{amsmath}
%\usepackage{amsfonts}
\usepackage{amssymb}
\usepackage{latexsym}
\usepackage{dsfont}
\usepackage{textcomp}

%% bera mit euler
\usepackage{bera,amssymb,eufrak}
\usepackage[euler-digits,euler-hat-accent]{eulervm}

\setlength{\parskip}{0.5em plus 1em minus 0.5em}
\clubpenalty 2000
\widowpenalty 2000
%\setlength{\tolerance}{2000}
\setlength{\emergencystretch}{20pt}   % kann man evtl. auf 20 erhöhen
%\setlength{\hfuzz}{1pt}
\setlength{\hfuzz}{10pt}

% ------------------------------------------------------------------------------
% pdf creation
\usepackage[%
  pdftex,
  colorlinks=true,
  urlcolor=red,
  citecolor=red,
  linkcolor=red,
  bookmarks=true,
  breaklinks=true
]{hyperref}

% ------------------------------------------------------------------------------
% graphics
\usepackage{siunitx}
\sisetup{scientific-notation = true}%
\usepackage[usenames,dvipsnames]{xcolor}
\usepackage{graphicx}
\usepackage{subfigure}
\usepackage{tikz}
\usetikzlibrary{calc,trees,positioning,arrows,chains,shapes.geometric,%
    decorations.pathreplacing,decorations.pathmorphing,shapes,%
    matrix,shapes.symbols}
\usepackage{pgfplots}
\pgfplotsset{compat=newest}
\usetikzlibrary{pgfplots.statistics} % For Box plots
\usepgfplotslibrary{units} % Allows to enter the units nicely


% ------------------------------------------------------------------------------
% listings
\usepackage{listings}		% include listings
\newcommand{\code}{\lstinline}
\lstloadlanguages{C++}
\lstset{
  basicstyle=\ttfamily\footnotesize,
  keywordstyle=\bfseries,
  stringstyle=\textit,
  emphstyle=\textit,
  commentstyle=\textit,
  breaklines=true,
  breakatwhitespace=true,
  breakautoindent=true,
  morekeywords={},
  numbers=left,
  numberstyle=\tiny\color{black},
  stepnumber=1,
  numbersep=5pt,
  tabsize=2,
  showspaces=false,
  showstringspaces=false,
    xleftmargin=17pt, 			% left margin where linenumbers and frames do start
    framexleftmargin=17pt, 		% left margin of frame (which will be colored)
    framesep=0pt,
    belowcaptionskip=0pt,
    captionpos=b,
    framerule=1pt
}

\usepackage{scrhack}
\usepackage{float}
\floatstyle{plain}
\newfloat{listing}{tbhp}{myc}
\floatname{listing}{Listing}

% ------------------------------------------------------------------------------
% tables
\usepackage{multirow}
%\usepackage{ragged2e}
\usepackage{tabularx}
\usepackage{longtable}
\usepackage{booktabs}   % toprule, midrule in Tabellen
\newcolumntype{Y}{>{\raggedright\arraybackslash}X}
\usepackage{caption}

\usepackage{collcell}
\newcolumntype{L}{>{\collectcell\lstinline}l<{\endcollectcell}}
\newcolumntype{R}{>{\collectcell\lstinline}r<{\endcollectcell}}

\usepackage{bytefield}

%------------------------------------------------------------------------------
\usepackage{xspace}
\newcommand{\mythos}{\textsc{MyThOS}\xspace}
\def\cc{{C\nolinebreak[4]\hspace{-.05em}\raisebox{.4ex}{\tiny\bf ++}}}

\usepackage[textsize=tiny]{todonotes}

%\usepackage{titling}
\usepackage{enumerate}
\usepackage{epigraph}
\renewcommand{\epigraphsize}{\footnotesize}
\setlength{\epigraphwidth}{0.618\textwidth}
\usepackage{placeins}

\titlehead{
\begin{center}
  \noindent\rule{1.0\textwidth}{1pt}\\[1cm]
  \includegraphics[scale=0.5]{MyThOS_Logo}\\[0.5cm]
  \sffamily
  \small
  \begin{tabularx}{\textwidth}{lX}
    Förderkennzeichen: & 01IH130003 \\
    Vorhabensbezeichnung: & MyThOS \\
    & Modulares Betriebssystem für Massiv Parallele Anwendungen
  \end{tabularx}\\[0.5cm]
  \noindent\rule{1.0\textwidth}{1pt}
\end{center}
%\vspace{1cm}
}


\usepackage{scrlayer-scrpage}
\makeatletter
\title{\@myTitle}
\renewcommand*{\title}[1]{\def\@myTitle{#1}}
\setkomafont{pageheadfoot}{\sffamily\small}
\lohead{%
\@myTitle\\%
\headmark%
}
\cohead{}
\rohead{\includegraphics[height=1.5cm]{MyThOS_Logo}}
\makeatother
\pagestyle{scrheadings}

%\usepackage{showframe}

% ------------------------------------------------------------------------------
% layout
\usepackage{multicol}
\usepackage{marginnote}
\renewcommand*{\marginfont}{\footnotesize}
\typearea[current]{last}		% recalculate type area


\title{MyThOS D4.2.2 Developer Tools}
\author{Stefan Bonfert, Vladimir Nikolov, Robert Kuban, Randolf Rotta}

\hypersetup{
  pdftitle={MyThOS D4.2.2 Developer Tools},
  pdfsubject={MyThOS Deliverable},
  pdfauthor={Stefan Bonfert, Vladimir Nikolov, Robert Kuban, Randolf Rotta},
  pdfkeywords={MyThOS} % comma separated
}

\begin{document}
\selectlanguage{ngerman}
\maketitle

\begin{abstract}
Zur Untersützung der Entwickler werden eine Reihe von Softwareentwicklungswerkzeugen zur Verfügung gestellt. Dabei wird soweit möglich auf bestehende Werkzeuge zurückgegriffen und durch Anpassung und Erweiterungen eine Umsetzung realisiert.
\end{abstract}

\newpage
\tableofcontents
\newpage
% --- content ------------------------------------------------------------------

\section{Einleitung}
Dieses Dokument fasst den aktuellen Stand der Softwareentwicklungswerkzeuge 
zusammen, die bei
der Entwicklung des Kernels und bei der Entwicklung bzw. Anpassung der 
Anwendungen an die neue
Ausführungsumgebung verwendet werden. Darüber hinaus werden auch Vorgänge und 
Tools
beschrieben, welche bei der Auswertung der Anwendungsperformance verwendet 
wurden und auch
solche, die beim aktuellen Betrieb der Anwendungen für das Erstellen und für 
die Auswertung von
Simulationstests benötigt werden. Letztere werden auch für die Evaluation der
Anwendungsperformance mit der neuen Betriebssystem-Umgebung zum Einsatz kommen.
Dieses Dokument dient somit der Erstellung eines minimalen Werkzeugkastens, 
welcher für die
Entwicklung von Anwendungen und deren Betrieb bzw. Evaluation mit MyThOS 
bereitgestellt wird.
Die verwendeten Werkzeuge sind größtenteils frei verfügbar (z.B. open source), 
proprietäre
Produkte, wie der Intel C/C++ Compiler, die z.B. für eine 
Performance-Steigerung der Anwendungen
auf der Evaluationshardware verwendet werden können, werden ebenfalls 
diskutiert, obwohl diese
nicht zwingendermaßen für die Entwicklung notwendig sind.
In diesem Dokument sind lediglich diejenigen Werkzeuge beschrieben, die seit 
der Veröffentlichung von Deliverable D4.2.1 verändert oder zusätzlich in den 
Entwicklungsprozess aufgenommen wurden.
\section{Build-Tools}

In \mythos werden zusammengehörende Dateien, die eine eigenständige Einheit 
bilden zu Modulen zusammengefasst. Dadurch können sowohl Anforderungen von 
Modulen wie besipeilsweise Compiler-Unterstützung, als auch Unterschiede 
zwischen den Kernel-Varianten für unterschiedliche Platformen berücksichtigt 
werden. Zudem erleichtert die modulorientierte Struktur die Evaluation 
verschiedener Implementierungsvarianten, da einzelne Module einfach 
austauschbar sind.

Jedes Modul stellt eine Menge von Dateien bereit, und hängt von anderen Dateien 
ab. Solche Abhängigkeiten können beispielsweise Header-Dateien oder spezielle 
Implementierungsvarianten einzelner Objekte darstellen. Die meisten 
Abhängigkeiten müssen nicht manuell spezifiziert werden, sondern werden 
automatisiert aus den \texttt{\#include} Anweisungen im Quellcode extrahiert.

Um automatisiert alle für eine bestimmte Konfiguration benötigten Dateien 
zusammenzustellen wird eine Beschreibung der gewünschten Konfiguration 
benötigt. Diese Spezifikation legt fest, welche Module für eine 
bestimmte Ziel-Konfiguration benötigt werden, um ein lauffähiges Kernel-Image 
zu erstellen.\newline
Zur Auflösung der hier beschriebenen Abhängigkeiten wird eine Liste aller 
benötigten Datien und Module erstellt. Anschließend wird ein Modul für das 
Kernel-Abbild ausgewählt, falls es eines der benötigten Objekte bereitstellt.
Falls zwei Module die selbe Datei bereitstellen entsteht ein Konflikt, der 
vom Entwickler manuell aufgelöst werden muss, indem er eines der möglichen 
Module auswählt. Zur einfacheren Auswahl größerer Mengen von Modulen können 
Tags verwendet werden, die Teil der Modulbeschreibung sind. So sind 
beispielsweise alle Module, die auf die Xeon Phi Karte zugeschnitten sind, mit 
dem Tag \texttt{platform:knc} markiert.

\begin{lstlisting}[float, label=lst:module, caption=Ein Beispiel einer 
Modulbeschreibung.]
[module.boot-memory-multiboot]
requires = ["platform:multiboot"]
incfiles = ["boot/memory/Stage3Setup.h"]
kernelfiles = ["boot/memory/Stage3Setup.cc", "boot/memory/Stage3Setup.cc", 
"boot/memory/Stage3Setup-multiboot.cc"]

[module.boot-memory-gem5]
requires = ["platform:gem5"]
incfiles = ["boot/memory/Stage3Setup.h"]
kernelfiles = ["boot/memory/Stage3Setup.cc", "boot/memory/Stage3Setup.cc", 
"boot/memory/Stage3Setup-e820.cc"]

[module.boot-memory-knc]
requires = ["platform:knc"]
incfiles = ["boot/memory/Stage3Setup.h"]
kernelfiles = ["boot/memory/Stage3Setup.cc", "boot/memory/Stage3Setup.cc", 
"boot/memory/Stage3Setup-sfi.cc", "boot/memory/Stage3Setup-knc.cc"]
\end{lstlisting}

Jedes Modul befindet sich mit seinen Quelldateien und der Spezifikation in 
einem eigenen Ordner. Die Spezifikation wird dabei in einer oder mehreren 
\texttt{*.module}-Dateien festgehalten. Die Pfade zu den einzelnen Quelldateien 
sind darin relativ zum Pfad der Spezifikation angegeben. Der Pfad zum 
Verzeichnis, das alle vorhandenen Module beinhaltet, ist Teil der 
Konfigurations-Spezifikation. Listing \ref{lst:module} zeigt die 
Spezifikation eines Mouls, das die Datei \texttt{Stage3Setup.h} für 
unterschiedliche Plattformen bereitstellt. Anhand des 
platform-Tags wird die korrekte Implementierungs-Variante augewählt.

Um ein komplettes Kernel-Image zu erstellen liest das \texttt{mcconf}-Tool 
alle vorhandenen Modulbeschreibungen ein, wählt anhand einer 
Konfigurations-Spezifikation Module aus um vorhandene Abhängigkeiten aufzulösen 
und kopiert die ausgewählten Quell-Dateien in einen Zielordner.
Außerdem generiert das Tool ein \texttt{Makefile}, das Regeln für die 
Übersetzung der Quelldateien und das Linken des fertigen Kernel-Images enthält.
Anschließend kann das \texttt{make} Kommando verwendet werden, um das 
Kernel-Abbild zu erstellen. Für die Adaption des \texttt{Makefile}s an die 
Anforderungen einzelner Module können Modulbeschreibungen Teile des 
\texttt{Makefile} enthalten.

Für die Konfiguration des Build-Prozesses können separate Module verwendet 
werden, die keine Quelldateien enthalten. Damit können beispielsweise 
verschiedene je nach Zielplattform  unterschiedliche Compiler verwendet werden.


\section{Debugging}
Für die Fehlersuche im Kernel und in Anwendungsprogrammen kamen zunächst 
hauptsächlich klassische Debugger in Kombination mit Emulatoren wie Qemu und 
Bochs zum Einsatz. Um jedoch die Ausführung von parallelem Code auf realer 
Hardware wie dem Xeon Phi Coprozessor nachvollziehen zu können, sind 
zusätzliche Werkzeuge notwendig, da die Ausführung des Codes nich beliebig 
unterbrochen werden kann.

Zu diesem Zweck wurde ein Tracing Werkzeug entwickelt, das verteilt über alle 
Hardwarethreads Traces sammelt und diese aggregiert ausgeben kann. Hierfür wird 
auf jedem Hardwarethread separater Speicher verwendet, um Synchronisation beim 
Sammeln von Traces zu vermeiden, da dadurch die Ausführung selbst beeinflusst 
werden würde.

Zusätzlich wurde ein neus Werkzeug entwickelt, mit dem kausale Zusammenhänge 
zwischen Ereignissen auf verschiedenen Hardwarethreads extrahiert und 
anschließend für die Programmanalyse verwendet werden können. Solche Ereignisse 
können beispielsweise durch Nachrichtenaustausch hervorgerufene 
Funktionsaufrufe sein.

Durch dieses Werkzeug ist es möglich, die Ausführung von parallelem Code im 
Nachhinein genau zu unteruschen und wichtige Informationen über Parallelität 
und Abhängigkeiten zu extrahieren. Dadurch können Deadlocks oder 
Kommunikationsmuster deutlich einfacher untersucht werden.

% ------------------------------------------------------------------------------
% \bibliographystyle{alpha}
% \bibliography{literature}

\end{document}
